\documentclass[french]{article}

\usepackage[utf8]{inputenc}
\usepackage[T1]{fontenc}
\usepackage{lmodern}
\usepackage{geometry}
\usepackage{graphicx}
\usepackage{babel}
\usepackage{amsmath}
\usepackage{amsthm}
\usepackage{amssymb}
\usepackage{hyperref}
\usepackage{stmaryrd}
\usepackage{enumitem}
\usepackage[most]{tcolorbox}
\usepackage{dsfont}
\usepackage{multirow}
\usepackage{etoc}
\usepackage{titlesec}
\usepackage{esint}
\usepackage{cancel}
\usepackage{mathdots}
\usepackage{indentfirst}
\usepackage{lastpage}
\usepackage{fancyhdr}

\pagestyle{fancy}
\fancyhf{}
\renewcommand{\headrulewidth}{0pt}
\renewcommand{\footrulewidth}{0pt}
\fancyhead[L,R]{}
\fancyfoot[R]{\thepage \ /\ \pageref{LastPage}}

\fancypagestyle{plain}{\fancyhead{}\renewcommand{\headrule}{}}

\setlength{\parindent}{0pt}

\setlist[itemize]{label=$\bullet$}
\setlist{before=\leavevmode, leftmargin=*}

\renewcommand{\P}{\mathbb{P}}
\newcommand{\N}{\mathbb{N}}
\newcommand{\Z}{\mathbb{Z}}
\newcommand{\Q}{\mathbb{Q}}
\newcommand{\R}{\mathbb{R}}
\newcommand{\C}{\mathbb{C}}
\newcommand{\K}{\mathbb{K}}
\renewcommand{\L}{\mathbb{L}}
\newcommand{\U}{\mathbb{U}}
\newcommand{\st}{^*}
\newcommand{\tq}{\middle|}
\newcommand{\inv}{^{-1}}
\newcommand{\E}{\mathbb{E}}
\newcommand{\F}{\mathbb{F}}
\newcommand{\V}{\mathbb{V}}
\renewcommand{\ker}{\text{Ker}}
\newcommand{\im}{\text{Im}}
\newcommand{\card}{\text{Card}}
\newcommand{\Tr}{\text{Tr}}
\newcommand{\Vect}{\text{Vect}}
\newcommand{\rg}{\text{rg}}
\newcommand{\vertiii}[1]{{\left\vert\kern-0.25ex\left\vert\kern-0.25ex\left\vert #1 \right\vert\kern-0.25ex\right\vert\kern-0.25ex\right\vert}}
\renewcommand{\o}{\text{o}}
\renewcommand{\O}{\text{O}}
\renewcommand{\d}{\text{d}}
\newcommand{\Id}{\text{Id}}


\theoremstyle{plain}
\newtheorem*{ind}{Indication}

\theoremstyle{definition}

\newtheorem*{ex}{Exemple}
\newtheorem*{rmq}{Remarque}
\newtheorem*{notation}{Notation}
\newtheorem*{rappel}{Rappel}
\newtheorem*{terminologie}{Terminologie}
\newtheorem*{attention}{Attention}
\newtheorem*{conséquence}{Conséquence}

\theoremstyle{remark}
\newtheorem*{app}{Application}

\newtcolorbox[auto counter]{dfn}[1][]{
	enhanced,
	colback=white,
	colframe=black,
	sharp corners,
	boxrule=0.8pt,
	fonttitle=\bfseries,
	colbacktitle=white,
	coltitle=black,
	attach boxed title to top left={yshift=-1mm,xshift=-0.5mm},
	boxed title style={size=minimal, right=2pt, sharp corners, colframe=white},
	title={Définition \thetcbcounter \ifstrempty{#1}{}{ (#1)}},
	before skip=0.5em,
	after skip=0.5em
}

\newtcolorbox[auto counter]{exo}[1][]{
	enhanced,
	arc=1mm,
	colback=white,
	colframe=black,
	coltitle=black,
	fonttitle=\bfseries,
	boxrule=0.8pt,
	shadow={1.2pt}{-1.2pt}{0pt}{black!50},
	attach title to upper={\ },
	title={Exercice \thetcbcounter\ifblank{#1}{}{ #1}},
	before skip=0.5em,
	after skip=0.5em,
	left=2mm, right=2mm, top=1mm, bottom=1mm
}

\newtcolorbox{met}[1][]{
	enhanced,
	arc=1mm,
	colback=white,
	colframe=black,
	coltitle=black,
	fonttitle=\bfseries,
	boxrule=0.8pt,
	shadow={1.2pt}{-1.2pt}{0pt}{black!50},
	attach title to upper={\ },
	title={Point méthode \ifblank{#1}{}{#1}},
	before skip=1em,
	after skip=1em,
	left=2mm, right=2mm, top=1mm, bottom=1mm
}

\newcounter{thmcount}

\newtcolorbox[use counter=thmcount]{thm}[1][]{
	enhanced,
	colback=white,
	colframe=black,
	sharp corners,
	left skip=0mm,
	boxrule=1.3pt,
	fonttitle=\bfseries,
	colbacktitle=white,
	coltitle=black,
	attach boxed title to top left={yshift=-1mm,xshift=-0.5mm},
	boxed title style={size=minimal, right=2pt, sharp corners, colframe=white},
	title={Théorème \thethmcount \ifstrempty{#1}{}{ (#1)}},
	before skip=0.5em,
	after skip=0.5em
}

\newtcolorbox[use counter=thmcount]{prop}[1][]{
	enhanced,
	colback=white,
	colframe=black,
	sharp corners,
	boxrule=0.8pt,
	fonttitle=\bfseries,
	colbacktitle=white,
	coltitle=black,
	attach boxed title to top left={yshift=-1mm,xshift=-0.5mm},
	boxed title style={size=minimal, right=2pt, sharp corners, colframe=white},
	title={Proposition \thethmcount \ifstrempty{#1}{}{ (#1)}},
	before skip=0.5em,
	after skip=0.5em
}

\newtcolorbox[use counter=thmcount]{cor}[1][]{
	enhanced,
	colback=white,
	colframe=black,
	sharp corners,
	boxrule=0.8pt,
	fonttitle=\bfseries,
	colbacktitle=white,
	coltitle=black,
	attach boxed title to top left={yshift=-1mm,xshift=-0.5mm},
	boxed title style={size=minimal, right=2pt, sharp corners, colframe=white},
	title={Corollaire \thethmcount \ifstrempty{#1}{}{ (#1)}},
	before skip=0.5em,
	after skip=0.5em
}

\newtcolorbox[use counter=thmcount]{lem}[1][]{
	enhanced,
	colback=white,
	colframe=black,
	sharp corners,
	boxrule=0.5pt,
	fonttitle=\bfseries,
	colbacktitle=white,
	coltitle=black,
	attach boxed title to top left={yshift=-1mm,xshift=-0.5mm},
	boxed title style={size=minimal, right=2pt, sharp corners, colframe=white},
	title={Lemme \thethmcount \ifstrempty{#1}{}{ (#1)}},
	before skip=0.5em,
	after skip=0.5em
}

\addto\captionsfrench{\renewcommand*{\contentsname}{Dénombrabilité}}

\title{Dénombrabilité}
\date{}

\begin{document}
\tableofcontents

\newpage
\maketitle

\section{Équipotence}

\subsection{Relation d'équipotence}

\begin{dfn}
	Deux ensembles sont \textbf{équipotents} s'il existe une bijection entre eux.
\end{dfn}
Il s'agit d'une relation d'équivalence. Sur quoi ?

\begin{ex}
	Si $E$ est un ensemble, alors $\mathcal{P}(E)$ et $\{0,1\}^E$ sont équipotents.
\end{ex}

\begin{prop}[Théorème de Cantor]%prop
	Soit $E$ un ensemble. Il n'existe pas de surjection de $E$ sur $\mathcal{P}(E)$.
\end{prop}
\begin{proof}
	Si $f$ est une application de $E$ dans $\mathcal{P}(E)$, on montre que l'ensemble $$\{x\in E\ : \ x\notin f(x)\}$$ n'a pas d'antécédent par $f$.
\end{proof}

\subsection{Parties de $\N$}

\begin{dfn}
	Un ensemble $E$ est \textbf{fini} s'il est équipotent à une partie de $\N$ de la forme $\llbracket1,n\rrbracket$. Dans ce cas, un tel entier $n$ est unique : on l'appelle \textbf{cardinal} de $E$ et on le note $\card(E)$.
\end{dfn}
\begin{proof}
	On montre par récurrence \begin{center}
		"s'il existe une bijection $f:\llbracket1,n\rrbracket\to\llbracket1,m\rrbracket$, alors $n=m$"
	\end{center} en se ramenant au cas où $f(n)=m$.
\end{proof}

\begin{prop}%prop
	Une partie $F$ d'un ensemble $E$ fini est elle-même finie, et l'on a $\card(F)\leqslant\card(E)$. On a de plus $F=E$ si, et seulement si, $\card(F)=\card(E)$.
\end{prop}
\begin{proof}
	On montre que si $a\in E$, alors $\card(E\setminus\{a\})=\card(E)-1$, puis on raisonne par récurrence sur $\card(E)$.
\end{proof}

\begin{prop}%prop
	Une partie de $\N$ est finie si, et seulement si, elle est majorée.
\end{prop}

\begin{rmq}
	En particulier, $\N$ n'est pas un ensemble fini.
\end{rmq}

\begin{prop}%prop
	Toute partie infinie de $\N$ est en bijection avec $\N$.
\end{prop}
\begin{proof}
	Si $A$ est une partie infinie de $\N$, on construit par récurrence la suite $(a_n)_{n\in\N}$ définie par $a_0=\min(A)$ et : $$\forall n\in\N,\ a_{n+1}=\min\{x\in A\ : \ x>a_n\}$$ et l'on vérifier que $n\mapsto a_n$ est une bijection de $\N$ sur $A$. \\
	Pour la surjectivité, on peut remarquer que $]a_{n},a_{n+1}[\cap A=\emptyset$.
\end{proof}

\section{Ensembles dénombrables}

\subsection{Définitions, exemples}

\begin{dfn}
	Un ensemble est \textbf{dénombrable} s'il est équipotent à $\N$. Il est \textbf{au plus dénombrable} s'il est équipotent à une partie de $\N$.
\end{dfn}

\begin{ex}
	\begin{enumerate}
		\item D'après la proposition précédente, les parties de $\N$ sont finies ou dénombrables.
		\item Un ensemble est au plus dénombrable si, et seulement s'il est fini ou dénombrable.
		\item $\Z$ est dénombrable.
	\end{enumerate}
\end{ex}

\begin{prop}%prop
	Un ensemble $A$ est au plus dénombrable si, et seulement s'il existe une injection de $A$ dans $\N$.
\end{prop}
\begin{proof}
	C'est la traduction de la définition.
\end{proof}

\begin{cor}%cor
	Toute partie d'un ensemble au plus dénombrable est au plus dénombrable.
\end{cor}

\begin{ex}
	L'ensemble $\N^2$ est dénombrable. En effet, il est infini et l'application $(p,q)\mapsto 2^p3^q$ est injective de $\N^2$ dans $\N$ par unicité de la décomposition en facteurs premiers.
\end{ex}

\begin{exo}
	Montrer que l'ensemble $\N^{(\N)}$ des suites presque nulles d'entiers naturels est au plus dénombrable.
\end{exo}

\begin{lem}%lem
	Soit $f:E\to F$ une application, avec $E$ non vide.
	\begin{itemize}
		\item $f$ est injective si, et seulement s'il existe $g:F\to E$ telle que $g\circ f=\text{Id}_E$.
		\item $f$ est surjective si, et seulement s'il existe $g:F\to E$ telle que $f\circ g=\text{Id}_E$.
	\end{itemize}
	Une telle application $g$ est alors respectivement surjective et injective.
\end{lem}%axiome du choix

\begin{prop}%prop
	Soit $A$ un ensemble non vide. Les propriétés suivantes sont équivalentes :
	\begin{enumerate}[label=(\roman*)]
		\item $A$ est au plus dénombrable ;
		\item il existe une surjection d'un ensemble au plus dénombrable $X$ sur $A$ ;
		\item il existe une surjection de $\N$ sur $A$.
	\end{enumerate}
\end{prop}
\begin{proof}
	Pour $(iii)\implies(i)$, on prend une surjection $f$ de $\N$ sur $A$ et l'on montre que l'application $a\mapsto\min f\inv(\{a\})$ est une injection de $A$ dans $\N$.
\end{proof}

\begin{met}
	Pour montrer qu'un ensemble infini est dénombrable, on peut :
	\begin{itemize}
		\item exhiber une bijection de $\N$ sur $A$, ce qui revient à numéroter les éléments de $A$ en une suite $(a_n)$ ;
		\item exhiber une injection de $A$ dans $\N$ (ou dans un ensemble dénombrable) ;
		\item exhiber une surjection de $\N$ (ou d'un ensemble dénombrable) dans $A$.
	\end{itemize}
\end{met}

\begin{exo}
	Montrer que $\N^2$ est dénombrable, en mettant en évidence une numérotation sur un dessin, puis en explicitant celle-ci.
\end{exo}

\begin{prop}%prop
	Soit $A$ un ensemble. Les propriétés suivantes sont équivalentes :
	\begin{enumerate}[label=(\roman*)]
		\item $A$ est infini ;
		\item il existe une suite d'éléments de $A$ distincts deux à deux ;
		\item $A$ contient une partie dénombrable.
	\end{enumerate}
\end{prop}

\begin{rmq}
	Un ensemble fini n'est donc pas dénombrable.
\end{rmq}

\subsection{Opérations sur les ensembles dénombrables}

\begin{exo}
	Montrer que la réunion de deux ensembles dénombrables est dénombrable. Et pour une réunion d'une famille finie non vide d'ensembles dénombrables ?
\end{exo}

\begin{prop}%prop
	Le produit cartésien de deux ensembles dénombrables est un ensemble dénombrable.
\end{prop}
\begin{proof}
	Soit $A$ et $B$ deux ensembles dénombrables ainsi que $f$, une bijection de $A$ dans $\N$, et $g$, une bijection de $B$ dans $\N$. \\
	Alors $h:(x,y)\mapsto(f(x),g(y))$ est une bijection de $A\times B$ dans $\N^2$, de réciproque $(p,q)\mapsto(f\inv(p),g\inv(q))$. Comme $\N^2$ est dénombrable, on en déduit que $A\times B$ est dénombrable.
\end{proof}

\begin{ex}
	$\Q$ est dénombrable.
\end{ex}

\begin{exo}
	Que peut-on dire du produit cartésien de deux ensembles au plus dénombrables ? A quelle condition est-il dénombrable ?
\end{exo}

\begin{cor}%cor
	Le produit cartésien d'une famille finie non vide d'ensembles dénombrables est dénombrable.
\end{cor}

\begin{exo}
	Montrer qu'un espace vectoriel sur $\Q$ de dimension finie non nulle est dénombrable.
\end{exo}

\begin{cor}%cor
	La réunion d'une famille au plus dénombrable d'ensembles au plus dénombrables est au plus dénombrable.
\end{cor}
\begin{proof}
	Soit $(A_i)_{i\in I}$ une famille d'ensembles non vides au plus dénombrables indexée par un ensemble $I$ non vide au plus dénombrable. En considérant, pour tout $i\in I$, une surjection $f_i$ de $\N$ sur $A_i$, on construit une surjection de $f:I\times\N$ sur $\displaystyle\bigcup_{i\in I}A_i$ en posant : $$\forall (i,n)\in I\times\N,\ f(i,n)=f_i(n)$$
\end{proof}

\begin{ex}
	On retrouve ainsi les résultats de l'exercice précédent.
\end{ex}

\begin{exo}
	Montrer que $\Q[X]$ est dénombrable. En déduire que l'ensemble des nombres algébriques, c'est-à-dire des éléments de $\C$ qui sont racine d'un polynôme non nul de $\Q[X]$, est dénombrable.
\end{exo}

\begin{prop}%prop
	Si $(a_i)_{i\in I}$ est une famille sommable de nombres complexes, l'ensemble $\{i\in I : a_i\ne 0\}$, appelé \textbf{support} de la famille $(a_i)_{i\in I}$, est au plus dénombrable.
\end{prop}

\begin{ex}
	L'ensemble des points de discontinuité d'une fonction monotone de $\R$ dans $\R$ est au plus dénombrable.
\end{ex}

\section{Exemples d'ensembles infinis non dénombrables}

Le corps $\Q$ des rationnels est dénombrable, comme on l'a vu précédemment. Le fait que $\R$ vérifie la propriété de la borne supérieure fait qu'il ne peut être dénombrable.

\begin{lem}%lem
	Soit $(u_n)_{n\in\N}$ une suite réelle.\\
	Il existe au moins un réel n'appartenant pas à $\{u_n \mid  n\in\N\}$.
\end{lem}
\begin{proof}
	On construit par récurrence deux suites adjacentes $(a_n)_{n\in\N}$ et $(b_n)_{n\in\N}$ vérifiant $\forall n\in\N,\  u_n\notin[a_n,b_n]$. Leur limite commune ne peut être égale à aucun des $u_n$.
\end{proof}

\begin{cor}%cor
	Le corps $\R$ est non dénombrable.
\end{cor}

Le théorème de Cantor prouve que $\mathcal{P}(\N)$ est indénombrable. On peut montrer que $\R$ et $\mathcal{P}(\N)$ sont équipotents.

\begin{exo}
	Montrer que si $A$ est un ensemble possédant au moins deux éléments, alors $A^\N$ est non dénombrable.
\end{exo}

\begin{exo}
	En utilisant le développement décimal illimité des réels, exhiber une injection de $\{0,1\}^\N$ dans $\R$. Exhiber une injection de $\R$ dans $\{0,1\}^\N$.
\end{exo}

\begin{rmq}
	Le théorème de Cantor-Bernstein stipule que s'il existe une injection d'un ensemble $E$ dans un ensemble $F$, et une injection de $F$ dans $E$, alors $E$ et $F$ sont équipotents. L'exercice précédent permet donc de montrer que $\R$ et $\mathcal{P}(\N)$ sont équipotents, ainsi donc que $\{0,1\}^\N$.
\end{rmq}

\section{Exercices}

\begin{exo}
	\begin{enumerate}
		\item Montrer qu'un ensemble est infini si, et seulement s'il contient une partie dénombrable.
		\item Montrer que si $A$ est un ensemble au plus dénombrable et $B$ un ensemble infini, alors $A\cup B$ est en bijection avec $B$.
		\item Montrer qu'il existe une bijection de $\R$ sur $\R_+$. Peut-on en trouver une continue ?
	\end{enumerate}
\end{exo}

\begin{exo}
	Soit $E$ un espace vectoriel contenant une famille génératrice dénombrable $(e_n)_{n\in\N}$. Montrer que toute famille libre de $E$ est au plus dénombrable.
\end{exo}

\begin{exo}
	Soit $P:\R^2\to\R$. On suppose que pour tout $(a,b)\in\R^2$, les fonctions partielles $x\mapsto P(x,b)$ et $y\mapsto P(a,y)$ sont polynomiales. Montrer que $P$ est polynomiale. \\
	Indication : pour $a\in\R$, on pourra considérer le degré de $y\mapsto P(a,y)$. \\
	Question subsidiaire : que se passe-t-il si l'on remplace $\R$ par $\Q$ ?
\end{exo}

\begin{exo}
	Une fonction polynomiale de $\R^2$ dans $\R$ est une combinaison linéaire de fonctions de la forme $(x,y)\mapsto x^ky^l$ avec $(k,l)\in\N^2$.\\
	Soit $(P_n)_{n\in\N}$ une suite de telles fonctions, supposées toutes non nulles. Montrer qu'il existe un couple $(a,b)\in\R^2$ tel que $\forall n\in\N,\ P_n(a,b)\ne 0$.\\
	La propriété subsiste-t-elle si l'on remplace partout $\R$ par $\Q$ ? \\
	Indication : on pourra écrire toute fonction polynomiale sous la forme $(x,y)\mapsto\displaystyle\sum_{l=0}^{N}Q_l(x)y^l$ avec $Q_0,\ldots,Q_N$ des fonctions polynomiales d'une seule variable.
\end{exo}

\end{document}